\chapter{Normative and Informative References}
\section{Normative Internal References}
\meta{State the internal documents whose requirements are imported or enforced.}{Controlled internal specifications.}{Normative reference list.}

The following are normative to this specification:
\begin{itemize}
    \item \textbf{UF Specification v1.4.0 Skeleton}: canonical L0--L4 semantics, invariants, stability metrics, financial adapter constraints, governance, and auditability.
    \item \textbf{UF-Core Kernel Security Specification with SES}: kernel security boundary, UF-SP handling, SES redaction/encryption requirements, DomainParams, envelope protection, deployment profiles, and the Threat \& Response Matrix.
    \item \texttt{TFE\_Specification\_v2\_5.tex}: controlled mathematical and governance specification baseline.
    \item \texttt{TFE\_5\_3\_Implementation Plan v2.6.md}: implementation plan.
    \item \texttt{PRODUCTION\_EVALUATION\_CONTRACT.md}: production evaluation contract.
    \item \texttt{TFE\_REARCHITECTURE\_MASTER\_PLAN\_v2\_7.md}: re-architecture master plan.
\end{itemize}

\section{Informative Internal References}
\meta{State the internal concept sources that inform the scientific design without overriding the normative sources.}{Project concept materials.}{Informative reference list.}

The following are informative to this specification:
\begin{itemize}
    \item \textbf{Information Resonance Field concepts}: mode/amplitude reasoning over histories, mixed-time structure, and negative-space semantics.
    \item \textbf{Deterministic Multi-View Ensemble (DMVE)} concepts: deterministic multi-view fusion, bounded diversity, and explicit verification style.
    \item \textbf{Coherence Field} concepts: tempo, phase, and coherent-potential language for process-state descriptions.
    \item \textbf{Deterministic self-organizing mosaic dynamics}: reproducible geometry-of-data intuition, attractor/repeller interactions, and bounded structural inconsistency as a useful signal.
\end{itemize}

\section{Informative External Alignment References}
\meta{State the documentary and quality-alignment frameworks used to shape the document and record controls.}{Common standards and guidance families.}{Alignment statement.}

This specification is aligned in documentary structure with expectations common to requirements engineering and regulated computerized systems, including requirements specification practice, ALCOA+/GDocP data-integrity expectations, and controlled change/validation documentation. Such alignment is documentary and procedural; certification requires validated implementation and organizational controls.

Informative external references include:
\begin{enumerate}[leftmargin=1.2cm]
    \item ISO/IEC/IEEE 15288:2023, Systems and software engineering --- System life cycle processes
    \item ISO/IEC/IEEE 29148:2018, Systems and software engineering --- Life cycle processes --- Requirements engineering
    \item ISO/IEC 25010:2023, Systems and software engineering --- Product quality model
    \item ISO 9001:2015, Quality management systems --- Requirements
    \item A-Team Insight, \emph{The Top Low Latency Data Feed Providers}, 2023
    \item Supermicro, \emph{Building High-Performance Infrastructure for Real-Time Financial Analytics}
    \item AVP, \emph{The Future of Financial Advice: Technology's Role in the Market}
\end{enumerate}
